% 摘要正文与关键词。本文件由 \begin{abstract} 环境包裹，不要在此处写 \section。
% Every number below also appears in the body; the abstract introduces none.

本文围绕大语言模型训练中的数据质量、领域配比、规模与算力，依次建立质量评价与配比模型、标度律、
算力约束下的资源配置模型与能力前沿预测模型，并按数据来源与识别强度区分结论的证据等级。

针对问题一，将 22 个质量指标统一为“越高越好”的秩分，按六个打分体系平衡聚合得到域级质量评分
（arXiv 最高，为 0.6594；GitHub 最低，为 0.4206），并在记录与指标对两个层面度量质量冲突。以含两两
交互项的岭回归刻画 17 域配比与损失的关系，同规模留出集上 $R^2=0.725$；但在 60M 与 1B 上损失水平
不能迁移，故不给出 10B/70B 的配比外推。质量评分作为相对代理量单独报告，不进入配比模型。

针对问题二，以 Huber 损失与 L-BFGS-B 算法稳健估计经典标度律。诊断表明主拟合数据 B1 与已发表公式
逐点一致，拟合只证明生成公式可以恢复；在独立观测数据 B5 与剔除同族行的 B4 上，中位绝对相对误差为
7.52\% 与 8.07\%。包含质量的候选形式因预设的采纳条件未通过、且质量轴没有第一方映射而\textbf{不予采纳}，
其边际效用、弹性与质量–参数等价条件只作为该形式下的解析推论给出。

针对问题三，把训练、长文本注意力与质量提升三项开销合成不等式算力约束，在经典标度律的有效域内求
损失最小的配置，并以约束激活状态的改变定义结构性转移。质量取基线、上下文长度为 4096 时，
$10^{19}$ FLOPs 为内点最优，$10^{22}$ FLOPs 时数据量达到有效域上界，$10^{24}$ FLOPs 时最优解为有效域
上角点，算力利用率仅 0.0245。5 个观测上下文长度下都恰有两个结构性转移点，上下文越长，转移点对应的
预算越高。由于质量基线没有数值、已采纳的标度律不含质量项，非基线的质量配置无法识别；三族原始质量
成本中没有一族在整个定义域上处处最低。以上配置结论均以经典标度律的有效域为前提。

针对问题四，按预先固定的规则冻结 1,385 个开源模型。损失–得分映射的可比锚点全部处于得分下限，
识别不出斜率，故直接以历史得分建模。以参数量为尺度，对话与微调模型平均得分的年增幅中约 10.9\%
与参数规模相关；在 54 个预训练模型子集上，约 65.7\% 与训练算力相关。以条件第 90 百分位数定义前沿，
12 个月延续基线为 50.89（含回测误差的 90\% 预测区间为 $[40.27, 61.52]$），24 个月压力外推为 60.80；
在给定转移假设下，算力增长减半与停止时 12 个月前沿分别为 47.64 与 44.39，属敏感性结果，而非已识别的
因果效应。

综上，四问沿同一证据链推进，肯定性与否定性结果共同界定了结论的适用范围。

\keywords{数据质量评价\quad 标度律\quad 算力约束优化\quad 结构性转移\quad 能力前沿预测}
